%% FullCourse_10Lecs -- Lecture 6, Topic 6.5a: Mean-Field Games (HJB side)
%% Source: ./archive_pre_split/Lec06_HA_Models.tex (section 10, first half)
%% Sub-split: original T5 (sections 10-11) exceeded 30-page ceiling at 31 pages.
%% Split into T5a (MFG motivation + HJB) and T5b (KFE + deep MFG + lab preview).
%% FullCourse-only content: Mean-Field Games section (not present in MiniCourse).

%% Shared preamble for the FullCourse_10Lecs topic decks.
%%
%% Each topic .tex \input's this file, then defines its own \title, \date
%% override (if any), \graphicspath, and \begin{document}.
%%
%% Reconstructed verbatim from the inlined copy preserved in
%% Lec10_Case_Studies/Lec10_T8_Case6_DeepRamsey.tex (lines 23-190), which is
%% explicitly annotated there as "from ../shared_preamble.tex, copied verbatim".
%% Base preamble matches MiniCourse_8hr/Slides/shared_preamble.tex; this version
%% additionally provides \sectiondivider, the conceptbox/cautionbox/examplebox/
%% takeawaybox tcolorboxes, and \compacttables used across the split decks.

\documentclass[10pt,english,aspectratio=169]{beamer}

\usepackage{ae,aecompl}
\usepackage[T1]{fontenc}
\usepackage{textcomp}
\usepackage[utf8]{inputenc}
\usepackage{babel}
\usepackage{datetime2}

\setcounter{secnumdepth}{3}
\setcounter{tocdepth}{3}

\usepackage{amsmath, amssymb, amsthm, graphicx}
\usepackage{booktabs, multirow}
\usepackage{tikz}
\usepackage{pgfplots}
\pgfplotsset{compat=1.16}
\usepackage[most]{tcolorbox}
\tcbuselibrary{skins, breakable, theorems}
\usepackage{listings}
\usepackage{xcolor}

\lstset{
  basicstyle=\ttfamily\small,
  keywordstyle=\color{blue!80!black}\bfseries,
  commentstyle=\color{green!50!black}\itshape,
  stringstyle=\color{orange!80!black},
  backgroundcolor=\color{gray!8},
  frame=single,
  rulecolor=\color{gray!40},
  breaklines=true,
  showstringspaces=false,
  numbers=none,
}

\lstdefinestyle{pythonstyle}{
    language=Python,
    basicstyle=\ttfamily\footnotesize,
    keywordstyle=\color{blue!80!black}\bfseries,
    commentstyle=\color{green!50!black}\itshape,
    stringstyle=\color{orange!80!black},
    frame=single,
    breaklines=true,
    showstringspaces=false,
    numbers=left,
    numbersep=5pt,
    numberstyle=\tiny\color{gray},
    backgroundcolor=\color{gray!8},
    rulecolor=\color{gray!40}
}

\ifx\hypersetup\undefined
  \AtBeginDocument{%
    \hypersetup{bookmarks=false,breaklinks=true,pdfborder={0 0 0},colorlinks=true,
      linkcolor=DarkText,urlcolor=RedTitle}%
  }
\else
  \hypersetup{bookmarks=false,breaklinks=true,pdfborder={0 0 0},colorlinks=true,
    linkcolor=DarkText,urlcolor=RedTitle}%
\fi

\makeatletter
\newcommand\makebeamertitle{%
  \begin{frame}[plain]
    \vspace*{1.0cm}
    \centering
    {\usebeamerfont{title}\usebeamercolor[fg]{title}\inserttitle\par}
    \vspace{1.0cm}
    {\usebeamerfont{author}\usebeamercolor[fg]{author}\insertauthor\par}
    \vspace{0.2cm}
    {\usebeamerfont{date}\usebeamercolor[fg]{date}\insertdate\par}
  \end{frame}}%
\AtBeginDocument{%
  \let\origtableofcontents=\tableofcontents
  \def\tableofcontents{\@ifnextchar[{\origtableofcontents}{\gobbletableofcontents}}
  \def\gobbletableofcontents#1{\origtableofcontents}%
}

\usetheme{metropolis}
\usefonttheme{professionalfonts}

\definecolor{LightGray}{RGB}{230,230,230}
\definecolor{RedTitle}{RGB}{0,0,200}
\definecolor{VeryLightGray}{RGB}{245,245,245}
\definecolor{DarkText}{RGB}{0,0,0}

\setbeamercolor{background canvas}{bg=white}
\setbeamercolor{title}{fg=RedTitle}
\setbeamercolor{author}{fg=DarkText}
\setbeamercolor{date}{fg=DarkText}
\setbeamercolor{frametitle}{bg=VeryLightGray,fg=DarkText}
\setbeamercolor{normal text}{fg=DarkText}
\setbeamerfont{title}{size=\large}

\setbeamersize{text margin left=5mm, text margin right=5mm}
\addtobeamertemplate{frametitle}{}{\vspace{-0.5em}}
\newcommand{\reducevspace}{\vspace{-0.5cm}}

% Shared section divider and box styles for split topic decks.
\newcommand{\sectiondivider}[2]{%
\begin{frame}[plain,noframenumbering]
\begin{center}
\vfill
{\Large\textbf{#1}\par}
\vspace{0.35cm}
{\normalsize #2\par}
\vfill
\end{center}
\end{frame}
}

\newtcolorbox{conceptbox}[1][]{
  colback=blue!3,
  colframe=RedTitle!55,
  boxrule=0.35mm,
  arc=1mm,
  left=2mm,
  right=2mm,
  top=1mm,
  bottom=1mm,
  #1
}
\newtcolorbox{cautionbox}[1][]{
  colback=red!3,
  colframe=red!50!black,
  boxrule=0.35mm,
  arc=1mm,
  left=2mm,
  right=2mm,
  top=1mm,
  bottom=1mm,
  #1
}
\newtcolorbox{examplebox}[1][]{
  colback=green!3,
  colframe=green!45!black,
  boxrule=0.35mm,
  arc=1mm,
  left=2mm,
  right=2mm,
  top=1mm,
  bottom=1mm,
  #1
}
\newtcolorbox{takeawaybox}[1][]{
  colback=VeryLightGray,
  colframe=RedTitle!55,
  boxrule=0.35mm,
  arc=1mm,
  left=2mm,
  right=2mm,
  top=1mm,
  bottom=1mm,
  #1
}
\newcommand{\compacttables}{\renewcommand{\arraystretch}{1.15}}

\setbeamertemplate{navigation symbols}{}
\setbeamertemplate{footline}{%
  \hfill%
  \usebeamercolor[fg]{page number in head/foot}%
  \usebeamerfont{page number in head/foot}%
  \insertframenumber\,/\,\inserttotalframenumber\kern1em\vskip2pt%
}

\makeatother

\author{Zhigang Feng}
% "date with time": compile timestamp so each build is identifiable.
\date{\today\ \DTMcurrenttime}


\graphicspath{{./pic/}{../../MiniCourse_8hr/Slides/Module2_DL_RL_Macro/pic/}{../}}

\title[AI for Econ Research]{AI for Economic Research: Dynamic Models, Language, and Agents\\
Lecture 6.5a: Mean-Field Games --- Motivation and HJB}

\begin{document}

\makebeamertitle

%=============================================================================
\begin{frame}{Topic 6.5a Agenda}
\begin{enumerate}
  \item \textbf{Mean-Field Games} --- motivation, history, two-equation system
  \medskip
  \item \textbf{HJB side} --- continuous-time control, Poisson income shocks
\end{enumerate}
\medskip
\textit{Arc: T1 Aiyagari + computation $\to$ T2a Why ML / Euler $\to$ T2b Global DNN $\to$ T3 KS / DeepHAM $\to$ T4 DEQN / OLG $\to$ \textbf{T5a MFG HJB} $\to$ T5b MFG KFE + Lab}
\end{frame}

%=============================================================================
\section{Mean-Field Games}
% FullCourse-only

\begin{frame}{Mean-Field Games: motivation}
\begin{itemize}
\item All the models we've seen so far are \emph{finite or large-but-finite} populations of agents (Aiyagari has a continuum but discretises numerically).
\medskip
\item \textbf{Mean-Field Games} (MFG) provide a continuous-time, continuum-of-agents formulation.
\medskip
\item Why care?
  \begin{itemize}
  \item Continuous time often \emph{simplifies} the math: PDEs replace functional equations.
  \item Powerful new toolbox for HA models with idiosyncratic and aggregate shocks.
  \item Tight link to optimal transport, stochastic analysis, and modern PDE theory.
  \item Recent breakthroughs in \emph{deep learning for MFG} (Carmona--Lauri\`ere, Achdou--Lasry--Lions, Han et al.).
  \end{itemize}
\medskip
\item This part is an introduction; you can think of MFG as the continuum limit of Aiyagari/KS.
\end{itemize}
\end{frame}

\begin{frame}{The core idea: continuum + mean-field feedback}
\begin{itemize}
\item Imagine a \emph{continuum} of identical agents, each solving an individual control problem.
\medskip
\item Each agent's payoff depends on the \emph{distribution} of all other agents (the ``mean field'').
\medskip
\item The distribution is itself determined by the optimal behaviour of the agents.
\medskip
\item This circular dependence is the defining feature of MFG --- and is exactly the same fixed point as in HA models, but stated in continuous time.
\medskip
\item Two mathematical objects describe the equilibrium:
  \begin{itemize}
  \item A backward Hamilton--Jacobi--Bellman (HJB) equation for the individual value function.
  \item A forward Kolmogorov forward / Fokker--Planck (KFE) equation for the distribution.
  \end{itemize}
\end{itemize}
\end{frame}

\begin{frame}{A brief history}
\begin{itemize}
\item \textbf{Engineering origins} (mid-2000s):\\
  Huang, Malham\'e \& Caines (2006, \emph{Communications in Information and Systems} 6(3):221--252) developed ``Nash certainty equivalence'' for large-population games.
\medskip
\item \textbf{Mathematical formalisation}:\\
  Lasry \& Lions (2007, \emph{Japanese Journal of Mathematics}) coined the term ``Mean-Field Games'' and developed the rigorous PDE theory.
\medskip
\item \textbf{Economic adoption}:\\
  Achdou--Han--Lasry--Lions--Moll (2014, 2022 \emph{REStud}) translated the framework into HA macroeconomics --- continuous-time Aiyagari--Bewley--Huggett.
\medskip
\item \textbf{Deep-learning era} (2020s):\\
  Carmona--Lauri\`ere, Han--Hu, Achdou--Pequito, $\ldots$ use deep neural networks to solve MFGs in high dimensions.
\end{itemize}
\end{frame}

\begin{frame}{The two-equation MFG system}
\textbf{Hamilton--Jacobi--Bellman (backward, individual)}:
\[
-\partial_t v(x,t) + \rho v(x,t) \;=\; H\!\big(x,\nabla v(x,t),m(\cdot,t)\big),
\]
where $H$ is the Hamiltonian of the individual control problem and $m(\cdot,t)$ is the population distribution at time $t$.
\medskip

\textbf{Kolmogorov forward / Fokker--Planck (forward, aggregate)}:
\[
\partial_t m(x,t) \;=\; -\nabla\!\cdot\!\big(s^*(x,t,m)\, m(x,t)\big) + \tfrac{1}{2}\Delta\!\big(\sigma^2(x,t)\,m(x,t)\big),
\]
where $s^*$ is the optimal drift derived from the HJB.
\medskip
\begin{itemize}
\item HJB and KFE are \emph{coupled}: the optimal control depends on $m$, and $m$ evolves under the optimal control.
\item In equilibrium, both equations must hold simultaneously.
\end{itemize}
\end{frame}

\begin{frame}{Why HJB is familiar, why KFE is new}
\begin{itemize}
\item \textbf{HJB} is the continuous-time limit of the Bellman equation. Macroeconomists have seen it many times: Merton's portfolio choice, deterministic optimal growth, $\ldots$
\medskip
\item \textbf{KFE / Fokker--Planck} describes the time evolution of a probability density under stochastic dynamics.
  \begin{itemize}
  \item Originated in physics: Fokker (1914), Planck (1917) --- diffusion of particles.
  \item Formalised by Kolmogorov (1931).
  \item Standard in physics and engineering, only recently mainstream in macro.
  \end{itemize}
\medskip
\item Together they describe the \emph{joint dynamics of optimal individual behaviour and the cross-sectional distribution} --- exactly what HA models need.
\end{itemize}
\end{frame}

\begin{frame}{HJB review: deterministic optimal control}
A generic optimal control problem:
\[
v(x_0) \;=\; \max_{\alpha(\cdot)} \int_0^\infty e^{-\rho t}\, r(x_t, \alpha_t)\, dt
\quad\text{s.t.}\quad \dot x_t = f(x_t, \alpha_t),\; x_0 \text{ given}.
\]
\medskip

The corresponding HJB equation:
\[
\rho\, v(x) \;=\; \max_{\alpha}\left\{r(x,\alpha) + v'(x)\,f(x,\alpha)\right\}.
\]
\medskip

\textbf{Interpretation}:
\begin{itemize}
\item LHS: ``required return'' on the value function.
\item RHS: instantaneous payoff plus the marginal value of the state moving in the chosen direction.
\item First-order condition over $\alpha$: $r_\alpha + v'(x) f_\alpha = 0$.
\end{itemize}
\end{frame}

\begin{frame}{From discrete-time Bellman to continuous-time HJB}
Start with the discrete-time Bellman over a small interval $\Delta$:
\[
v(x) \;\approx\; \max_{\alpha}\!\left\{r(x,\alpha)\,\Delta + e^{-\rho\Delta}\, v(x + f(x,\alpha)\Delta)\right\}.
\]
\medskip
Subtract $v(x)$ from both sides, divide by $\Delta$, and let $\Delta \to 0$ (using $e^{-\rho\Delta} \approx 1 - \rho\Delta$ and $v(x + f\Delta) \approx v(x) + v'(x) f\Delta$):
\[
0 \;=\; \max_\alpha\!\left\{r(x,\alpha) - \rho v(x) + v'(x) f(x,\alpha)\right\}.
\]
\medskip
Rearrange:
\[
\rho v(x) \;=\; \max_\alpha\!\left\{r(x,\alpha) + v'(x) f(x,\alpha)\right\}.
\]
\medskip
This is exactly the HJB.
\end{frame}

\begin{frame}{Stationary HJB and FOCs}
For the optimal growth problem with state $k$, control $c$, return $u(c)$, and transition $\dot k = f(k) - \delta k - c$:
\[
\rho\, v(k) \;=\; \max_c\!\left\{u(c) + v'(k)\,\big(f(k)-\delta k - c\big)\right\}.
\]
\medskip
First-order condition:
\[
u'(c) \;=\; v'(k) \quad\Longrightarrow\quad c^*(k) = (u')^{-1}(v'(k)).
\]
\medskip

Substituting back:
\[
\rho\, v(k) \;=\; u\!\big(c^*(k)\big) + v'(k)\,\big(f(k) - \delta k - c^*(k)\big).
\]
\medskip
A second-order ODE in $v$ to be solved on $k \in [0, \bar k]$ with appropriate boundary conditions.
\end{frame}

\begin{frame}{Poisson income shocks: state space}
\begin{itemize}
\item In the continuous-time Aiyagari, individual income $y$ jumps between two states $\{y_1, y_2\}$ at \emph{Poisson rates} $\lambda_1, \lambda_2$.
\medskip
\item State of an individual: $(a, y_i)$ with $i \in \{1,2\}$.
\medskip
\item Two value functions $v_1(a)$ and $v_2(a)$, one per income state.
\medskip
\item Asset evolves continuously: $\dot a = y_i + r a - c$.
\medskip
\item Income switches at Poisson rate; from state $i$, the agent moves to state $-i$ at rate $\lambda_i$.
\medskip
\item This setup is the natural continuous-time analogue of the discrete-time Aiyagari with Markov income.
\end{itemize}
\end{frame}

\begin{frame}{Deriving the HJB with jump terms \textsc{(advanced --- skip on first pass)}}
Over a short interval $\Delta$, with probability $\lambda_i \Delta$ the agent's income jumps from $y_i$ to $y_{-i}$. The Bellman is
\[
v_i(a) \approx \max_c\!\left\{u(c)\Delta + e^{-\rho\Delta}\!\left[(1-\lambda_i\Delta)\, v_i(a + \dot a\Delta) + \lambda_i\Delta\, v_{-i}(a)\right]\right\}.
\]
\medskip

Letting $\Delta \to 0$ and rearranging:
\[
\rho\, v_i(a) \;=\; \max_c\!\left\{u(c) + v_i'(a)\,(y_i + r a - c) + \lambda_i\,\big(v_{-i}(a) - v_i(a)\big)\right\}.
\]
\medskip
\begin{itemize}
\item Three terms: instantaneous utility, drift term (deterministic capital accumulation), jump term (Poisson income switch).
\item This is the HJB with jumps; the system is $i \in \{1,2\}$.
\end{itemize}
\end{frame}

\begin{frame}{Economic intuition of the jump term}
\[
\lambda_i \big(v_{-i}(a) - v_i(a)\big)
\]
\medskip
\begin{itemize}
\item This is the \emph{option value} of switching income state.
\medskip
\item If $v_{-i}(a) > v_i(a)$ (the other income state would yield higher continuation value at the same wealth), the jump term is \emph{positive} --- the agent enjoys an expected gain from the future income switch.
\medskip
\item Conversely, if $v_{-i}(a) < v_i(a)$, the jump term is a cost.
\medskip
\item The Poisson rate $\lambda_i$ controls the strength of this option value.
\medskip
\item Intuition: precautionary saving in continuous time appears partly through this jump term, in addition to the curvature of $u$ along the deterministic drift.
\end{itemize}
\end{frame}

\begin{frame}{Continuous-time Aiyagari: setup}
\textbf{Households.} A continuum of households indexed by current asset $a \in [\underline a, \infty)$ and income $y_i \in \{y_1, y_2\}$ with Poisson switching rates $\lambda_1, \lambda_2$.
\medskip

\textbf{Firm.} Representative firm with $Y = K^\alpha L^{1-\alpha}$ and competitive factor markets:
\[
r = \alpha (K/L)^{\alpha-1} - \delta, \qquad w = (1-\alpha) (K/L)^\alpha.
\]
\medskip

\textbf{Aggregation.} Aggregate capital and labour:
\[
K = \sum_{i=1}^{2} \int a\, g_i(a)\, da, \qquad
L = \sum_{i=1}^{2} y_i \int g_i(a)\, da,
\]
where $g_i(a)$ is the cross-sectional density in income state $i$.
\end{frame}

\begin{frame}{The aggregation loop}
\begin{enumerate}
\item Guess prices $(r, w)$ (or equivalently aggregate $K$).
\medskip
\item Solve the household HJB to get optimal consumption $c_i^*(a)$ and drift $s_i^*(a) = y_i + r a - c_i^*(a)$.
\medskip
\item Solve the KFE for the stationary density $g_i^*(a)$.
\medskip
\item Compute aggregate $K^* = \sum_i \int a\, g_i^*(a)\, da$.
\medskip
\item Update $(r,w)$ from the firm FOCs at $K^*$, and iterate until $K^* = K$.
\end{enumerate}
\medskip

This is structurally identical to the discrete-time Aiyagari nested fixed point, but every step is now an ODE/PDE rather than a discrete-time DP problem.
\medskip

\emph{Next (T5b):} the KFE side --- forward equation for the density, upwind discretisation, and the lab preview.
\end{frame}

\end{document}
